% !TeX root = ../main.tex
% !TeX encoding = utf8

\chapter{Vínculos laterales y centralización}
% Responsable: Ismael Sallami Moreno
% Secciones del índice:
%   4.3 Diseño de los vínculos laterales
%   4.4 Centralización y descentralización

\section{Diseño de los vínculos laterales}

\subsection{Planificación de acciones}

% ¿Existe algún plan de acción o campaña comercial que vincule a varias
% unidades o sucursales?
%   - Indicar qué unidades están implicadas.
%   - Describir cómo se comunica y gestiona el plan.

\subsection{Control del rendimiento}

% ¿Existen mecanismos de control del rendimiento?
%   - ¿Qué unidad fija los objetivos (sede central, dirección territorial...)?
%   - ¿Qué unidades están obligadas a cumplirlos?
%   - ¿Con qué frecuencia se revisan? ¿Qué ocurre si no se alcanzan?

\subsection{Dispositivos de enlace}

% ¿Qué dispositivos de enlace existen?
%   - Indicar el tipo (puesto de enlace, grupo de trabajo, director integrante, estructura matricial).
%   - ¿A qué personas o unidades relaciona?
%   - ¿Cuál es su objetivo?

\section{Centralización y descentralización}

% Indicar el tipo y grado de descentralización, tanto vertical como horizontal.
%
% Descentralización vertical:
%   - ¿Qué decisiones puede tomar la directora de sucursal de forma autónoma?
%   - ¿Qué decisiones requieren aprobación de la dirección territorial o sede central?
%
% Descentralización horizontal:
%   - ¿Los empleados de la sucursal tienen poder de decisión en su área de trabajo?
%   - ¿Existen especialistas o expertos con influencia informal en las decisiones?
%
% Especificar claramente a quién se delega poder y sobre qué tipo de decisiones.

\endinput
% -------------------------------------------------------------------
% FIN DEL CAPÍTULO — Ismael Sallami Moreno
% -------------------------------------------------------------------
